\documentclass[aspectratio=169,11pt]{beamer}
\usepackage[utf8]{inputenc}
\usepackage[T1]{fontenc}
\usepackage[default,scale=0.92]{sourcesanspro}
\renewcommand{\familydefault}{\sfdefault}
\usepackage[spanish,es-noshorthands]{babel}
\usepackage{fontawesome5}
\usepackage{booktabs}
\usepackage{amsmath}
\usepackage{tcolorbox}

% ── Paleta: generada desde estilo_infografia.PALETA (fuente única) ────────────
\definecolor{fondoOscuro}{HTML}{0D1B2A}%             Dark navy — fondo banda título
\definecolor{verdeTurquesa}{HTML}{004D40}%           Verde turquesa — bandas de marca
\definecolor{azulNoche}{HTML}{1B3A6B}%               Midnight blue — secciones, marcos
\definecolor{azulMedio}{HTML}{2A5298}%               Azul medio — variante de acento
\definecolor{cyanNeon}{HTML}{00C8E0}%                Cyan eléctrico — acentos primarios
\definecolor{verdeSignal}{HTML}{00B887}%             Verde señal — fortalezas, éxito
\definecolor{naranjaVivo}{HTML}{FF7043}%             Naranja vivo — mejoras, advertencias
\definecolor{rojoAlerta}{HTML}{E53935}%              Rojo alerta — errores
\definecolor{grisPapel}{HTML}{F4F6FA}%               Fondo tarjetas claras
\definecolor{grisLinea}{HTML}{C8D0E0}%               Bordes sutiles
\definecolor{grisTexto}{HTML}{6B7A99}%               Texto secundario
\definecolor{amarilloNota}{HTML}{FFB300}%            Notas / advertencias doradas
\definecolor{codigoFondo}{HTML}{1E2D3D}%             Fondo del bloque de código
\definecolor{codigoComentario}{HTML}{7A8FA6}%        Comentarios
\definecolor{codigoCadena}{HTML}{FF9D5C}%            Cadenas / strings
\definecolor{codigoKeyword}{HTML}{00C8E0}%           Palabras clave
\definecolor{codigoNumero}{HTML}{00E5A0}%            Números

\setbeamercolor{structure}{fg=azulMedio}
\setbeamercolor{title}{fg=fondoOscuro,bg=white}
\setbeamercolor{frametitle}{fg=white,bg=azulNoche}
\setbeamerfont{frametitle}{size=\large,series=\bfseries}
\setbeamercolor{block title}{fg=white,bg=azulMedio}
\setbeamercolor{block body}{bg=grisPapel,fg=fondoOscuro}
\setbeamercolor{block title alerted}{fg=white,bg=rojoAlerta}
\setbeamercolor{block body alerted}{bg=rojoAlerta!8,fg=fondoOscuro}
\setbeamercolor{block title example}{fg=white,bg=verdeTurquesa}
\setbeamercolor{block body example}{bg=verdeSignal!8,fg=fondoOscuro}
\setbeamertemplate{navigation symbols}{}
\setbeamertemplate{itemize item}{\color{cyanNeon}\faIcon{angle-right}~}
\setbeamertemplate{itemize subitem}{\color{azulMedio}\faIcon{circle}~}

\newtcolorbox{panelMono}[1][]{colback=codigoFondo,colframe=codigoFondo,boxrule=0pt,
  arc=3pt,coltext=white,fontupper=\ttfamily\scriptsize,
  left=7pt,right=7pt,top=6pt,bottom=6pt,#1}
\newtcolorbox{panelRegla}[1][]{colback=cyanNeon!10,colframe=cyanNeon,boxrule=1pt,
  arc=4pt,left=9pt,right=9pt,top=7pt,bottom=7pt,#1}

% Cifra grande + su leyenda. Se usa en minipage, así que nada de \par.
\newcommand{\cifra}[2]{%
  \begin{minipage}[t]{0.30\textwidth}\centering
  {\fontsize{28}{30}\selectfont\bfseries\color{azulNoche}#1}\\[3pt]
  {\footnotesize\color{grisTexto}#2}
  \end{minipage}}

\setbeamertemplate{footline}{%
  \leavevmode\hbox{%
  \begin{beamercolorbox}[wd=\paperwidth,ht=3ex,dp=1.6ex,leftskip=1em,rightskip=1em]
    {author in head/foot}
    \scriptsize\color{grisTexto}ELECTRONICA \enspace \textbar\enspace Patrón de arquitectura de 3 capas\hfill\insertframenumber/\inserttotalframenumber
  \end{beamercolorbox}}}

\title{IA en ingeniería con un esquema determinista de tres capas}
\author{ELECTRONICA}
\date{}

\begin{document}

\begin{frame}[plain]
  \begin{beamercolorbox}[sep=16pt,center]{title}
    {\footnotesize\bfseries\color{cyanNeon}\faIcon{cogs}\ INGENIERÍA DE DATOS}\\[12pt]
    {\LARGE\bfseries\color{fondoOscuro} IA en ingeniería con un\\ esquema determinista de tres capas}\\[12pt]
    {\large\color{azulMedio} Por qué el 84 \% de un sistema con LLM no debería contener un LLM}\\[20pt]
    {\footnotesize\color{grisTexto}\faIcon{user-circle}~ELECTRONICA \qquad
      \faIcon{calendar}~2026-10-05\qquad
      \faIcon{folder-open}~51 directivas, 62 scripts, 23 flujos}
  \end{beamercolorbox}
\end{frame}

\begin{frame}{Un LLM puede cumplir el contrato entero \emph{y aun así mentir}}
  \begin{alertblock}{Lo que devolvió el modelo}
    \begin{panelMono}
puntaje\_sugerido : "5.5/10"       \textcolor{codigoCadena}{<- inventado}\\
nivel\_desempeno  : "Bueno"        \textcolor{codigoCadena}{<- inventado}\\
observaciones    : [ "0.5/1", "0.7/1", "0.3/1" ]   \textcolor{codigoNumero}{<- suma 1.5}
    \end{panelMono}
  \end{alertblock}
  \smallskip
  \begin{block}{Obsérvese lo que \emph{sí} fue correcto}
    \begin{itemize}
      \item JSON válido \quad \item campos presentes \quad \item denominadores correctos
      \item \ldots y el total \textbf{no es la suma de sus propias partes}.
    \end{itemize}
  \end{block}
  \begin{panelRegla}
    Un LLM no es una calculadora. Cumplir un formato no es cumplir un contrato.
  \end{panelRegla}
\end{frame}

\begin{frame}{La asimetría que hay que resolver}
  \begin{columns}[T]
    \begin{column}{0.48\textwidth}
      \begin{alertblock}{Lo que es un LLM}
        \begin{itemize}
          \item Un generador de texto \emph{plausible}
          \item Sin garantía de aritmética
          \item Sin garantía de repetibilidad
          \item Y sin que eso sea un defecto: es su naturaleza
        \end{itemize}
      \end{alertblock}
    \end{column}
    \begin{column}{0.48\textwidth}
      \begin{exampleblock}{Lo que exige la ingeniería}
        \begin{itemize}
          \item El mismo input $\rightarrow$ el mismo output
          \item Que se pueda auditar \emph{por qué}
          \item Que el error sea ruidoso, no silencioso
        \end{itemize}
      \end{exampleblock}
    \end{column}
  \end{columns}
  \smallskip
  \begin{panelRegla}
    No es una discusión filosófica. Es una asimetría \textbf{medible}, y se mide
    en fallos reales con coste concreto.
  \end{panelRegla}
\end{frame}

\begin{frame}{La respuesta: separar por responsabilidad}
  \begin{center}
  \begin{tabular}{@{}c c c l@{}}
    \toprule
    \textbf{Capa} & \textbf{Artefacto} & \textbf{Responde} & \textbf{Contiene LLM} \\
    \midrule
    1 \quad Directiva & \texttt{directives/*.yaml} & ¿qué? & no \\
    2 \quad Orquestación & \texttt{flujo\_*.py}, \texttt{mcp\_*} & ¿cuándo? & no (solo decide) \\
    3 \quad Ejecución & \texttt{execution/*.py} & ¿cómo? & 10 de 62 \\
    \bottomrule
  \end{tabular}
  \end{center}
  \smallskip
  \begin{block}{La regla que sostiene el patrón}
    La lógica de negocio no se delega al modelo. Se le pide lo único que hace
    bien --\emph{producir texto--} y el resto se calcula en código.
  \end{block}
\end{frame}

\begin{frame}{La frontera, con números medidos}
  \vspace{2pt}
  \cifra{84 \%}{scripts de ejecución sin ningún LLM}
  \hfill
  \cifra{10}{fronteras con LLM}
  \hfill
  \cifra{\$0}{coste de decidir el modelo}
  \par\vspace{14pt}
  \begin{block}{El enrutador es una función pura}
    Elegir qué modelo usar \emph{no lo decide un modelo}: lo decide código local,
    a partir de un descriptor medido. Mismo descriptor $\rightarrow$ mismo tier,
    siempre. Y no hace ni una sola llamada a la API.
  \end{block}
  \begin{alertblock}{Por qué esto no es un detalle de implementación}
    Las 10 fronteras están \textbf{declaradas}: son ficheros
    concretos y nombrados, no un \texttt{if} repartido por el código.
  \end{alertblock}
\end{frame}

\begin{frame}{Capa 1 \textbullet\ La directiva: el \emph{qué}}
  \begin{panelMono}
goal: >
  Evaluar el examen del estudiante enviando las páginas del PDF
  como imágenes al modelo y obteniendo una evaluación en JSON.
required\_inputs:
  - name: pdf\_examen
    description: Ruta al PDF del examen.
steps:
  - step: 1
    description: Renderizar el PDF a 250 DPI y enviarlo en bloque
                  al modelo con la instrucción de corrección.
  \end{panelMono}
  \smallskip
  \begin{block}{Por qué está en YAML y no en el prompt}
    Porque es \textbf{configuración}, no conversación: se versiona, se revisa con
    un diff, y lo escribe alguien que entiende el dominio, no el modelo.
  \end{block}
\end{frame}

\begin{frame}{Capa 2 \textbullet\ La orquestación: el \emph{cuándo}}
  \begin{columns}[T]
    \begin{column}{0.5\textwidth}
      \begin{block}{Lo que sí hace}
        \begin{itemize}
          \item Leer la directiva
          \item Resolver entradas y decidir el orden
          \item Delegar cada paso en capa 3
          \item Validar salidas y códigos de retorno
          \item Registrar estado para poder reanudar
        \end{itemize}
      \end{block}
    \end{column}
    \begin{column}{0.5\textwidth}
      \begin{alertblock}{Lo que NO hace}
        \begin{itemize}
          \item \ldots calcular
          \item \ldots sumar, comparar ni convertir
          \item \ldots procesar datos crudos
          \item \ldots decidir un modelo a ojo
        \end{itemize}
      \end{alertblock}
    \end{column}
  \end{columns}
  \begin{panelRegla}
    Una capa de decisión delgada, con la lógica fina abajo.
    Por eso el mismo fallo no se propaga a veinte sitios.
  \end{panelRegla}
\end{frame}

\begin{frame}{Capa 3 \textbullet\ La ejecución: el \emph{cómo}}
  \begin{itemize}
    \item Funciones puras donde se puede: mismo input, mismo output, sin red.
    \item Cada script valida \emph{sus propias} salidas y falla ruidosamente.
    \item Códigos de salida tipados; un 0 significa 0, no \"no vi nada\".
    \item Cobertura de pruebas como forma de documentar el contrato.
  \end{itemize}
  \smallskip
  \begin{exampleblock}{Dónde vive la verdad}
    Cuando hay una discrepancia, la respuesta no la da el prompt ni lo que el
    modelo recuerda: la da el test que falla.
  \end{exampleblock}
\end{frame}

\begin{frame}{Fallo 1 \textbullet\ El modelo no suma: el programa suma}
  \begin{columns}[T]
    \begin{column}{0.48\textwidth}
      \begin{alertblock}{Síntoma}
        Salida con \texttt{status: ok}, JSON perfecto, cada ítem con su
        denominador \ldots y un total que no era la suma.
      \end{alertblock}
      \begin{block}{Causa profunda}
        Se le pidió al modelo el \emph{total}. La aritmética la hacía el LLM, y
        la aritmética es justo lo que un modelo no garantiza.
      \end{block}
    \end{column}
    \begin{column}{0.48\textwidth}
      \begin{exampleblock}{Arreglo}
        \begin{itemize}
          \item El esquema \emph{deja de pedir} el total
          \item El programa suma: función pura, testeable, sin red
          \item El máximo es un invariante, no una consecuencia
          \item Si se pasa del tope, se recorta \emph{y se avisa}
        \end{itemize}
      \end{exampleblock}
    \end{column}
  \end{columns}
  \begin{panelRegla}
    Lo que se transfiere: en cuanto un \textbf{número importa}, el número lo
    calcula el programa.
  \end{panelRegla}
\end{frame}

\begin{frame}{Fallo 2 \textbullet\ La escala la declara el instrumento}
  \begin{columns}[T]
    \begin{column}{0.47\textwidth}
      \begin{alertblock}{El síntoma}
        3 ítems de 1 punto (total 3) con la mitad respondidos. Al sumar en crudo:
        \textbf{1.5/10} y nivel \emph{Insuficiente}. El alumno sacaba el 50 \%.
      \end{alertblock}
      \begin{alertblock}{Peor: el bug espejo}
        Si el modelo se comía un criterio, su peso desaparecía del denominador
        y el resto se normalizaba \textbf{hacia arriba}: un fallo del modelo
        bonificaba al alumno.
      \end{alertblock}
    \end{column}
    \begin{column}{0.47\textwidth}
      \begin{exampleblock}{El arreglo}
        \begin{itemize}
          \item La escala la fija el \emph{instrumento} (la rúbrica), no quien
                lo corrige
          \item Normalizar a la escala institucional, siempre
          \item Un criterio ausente vale 0 sobre esa escala: no se renormaliza
          \item Sin escala declarada \emph{no se publica nota}, y se dice por qué
        \end{itemize}
      \end{exampleblock}
    \end{column}
  \end{columns}
\end{frame}

\begin{frame}{Fallo 3 \textbullet\ El modo JSON no es un esquema}
  \begin{panelMono}
request:  response\_format = \{"type": "json\_object"\}   \textcolor{codigoCadena}{<- solo promete un objeto}\\
respuesta: \{"text": "...", "preguntas": []\}          \textcolor{codigoCadena}{<- objeto válido. otro esquema.}
    \end{panelMono}
  \smallskip
  \begin{block}{El agravante}
    \texttt{response\_format} promete \textbf{una sola cosa}: que la respuesta
    sea un objeto. No promete las claves, ni los tipos, ni el vocabulario. Y el
    esquema del prompt es texto: la capa más débil del contrato.
  \end{block}
  \begin{exampleblock}{El arreglo}
    La barrera real es un \textbf{validador en el programa}: vocabulario cerrado
    para los campos enumerados, comprobación de tipos, y la clave que decide la
    vialidad del flujo es \emph{obligatoria}. Salida con código 5 y la lista de
    problemas; el JSON del modelo no se acepta como informe.
  \end{exampleblock}
\end{frame}

\begin{frame}{Los tres fallos, una sola regla}
  \begin{panelRegla}
    \centering\Large\bfseries
    Cuando un número importa,\\ el número lo calcula el programa.
    \par\vspace{8pt}
    \normalsize\normalfont
    Y cuando la forma importa, la forma la verifica el programa.
  \end{panelRegla}
  \vspace{8pt}
  \begin{columns}[T]
    \begin{column}{0.47\textwidth}
      \begin{block}{Aplicaciones fuera de la IA}
        \begin{itemize}
          \item Metrología y conversión de unidades
          \item Control de calidad y criterios de aceptación
          \item Presupuestos y dosages
          \item Cualquier cosa donde un decimal de más sea un problema
        \end{itemize}
      \end{block}
    \end{column}
    \begin{column}{0.47\textwidth}
      \begin{block}{Por qué así se aprende}
        \begin{itemize}
          \item No es desconfiar del modelo: es asignar la aritmética a lo único
                que sabe hacer aritmética.
        \end{itemize}
      \end{block}
    \end{column}
  \end{columns}
\end{frame}

\begin{frame}{Una idea formal que se transfiere: la álgebra de veredictos}
  \begin{columns}[T]
    \begin{column}{0.47\textwidth}
      \begin{alertblock}{La regla}
        \begin{itemize}
          \item Gana el \textbf{peor} estado, nunca el promedio
          \item Lo no verificado \textbf{impide} el verde
          \item Se distingue \emph{no instalado} de \emph{caído} de \emph{parado}
        \end{itemize}
      \end{alertblock}
    \end{column}
    \begin{column}{0.47\textwidth}
      \begin{exampleblock}{Por qué importa fuera de aquí}
        \begin{itemize}
          \item Es un promedio con tres estados en rojo y uno en amarillo: se
                ve verde mientras algo crítico no se midió
          \item Es exactamente la diferencia entre \emph{sin datos} y
                \emph{con buenos datos}
          \item Es control robusto aplicado a un pipeline de verificación
        \end{itemize}
      \end{exampleblock}
    \end{column}
  \end{columns}
  \begin{panelRegla}
    Y hay una segunda: \textbf{una sola definición}, compartida por las dos
    auditorías. Dos álgebras que divergen en silencio son la forma más cara de
    equivocarse.
  \end{panelRegla}
\end{frame}

\begin{frame}{Los límites, que también son parte del patrón}
  \begin{columns}[T]
    \begin{column}{0.47\textwidth}
      \begin{alertblock}{Lo que el patrón \emph{no} resuelve}
        \begin{itemize}
          \item La capa 1 (directivas) aún no es verificable por auditoría:
                es informativa
          \item Un motor de asistente rotativo no es reproducible
          \item El hardware disponible no entrena modelos
          \item Sin red, parte de los proveedores no son accesibles
        \end{itemize}
      \end{alertblock}
    \end{column}
    \begin{column}{0.47\textwidth}
      \begin{exampleblock}{Cómo se presenta esto}
        \begin{itemize}
          \item Un patrón cuya frontera se declara se enseña mejor que uno que
                se vende
          \item La reproducibilidad tampoco es gratis: cuesta discipline
                de capa 3 y tests
        \end{itemize}
      \end{exampleblock}
    \end{column}
  \end{columns}
\end{frame}

\begin{frame}{Cómo adoptarlo en su propio proyecto}
  \begin{enumerate}
    \item \textbf{Escriba el SOP antes del código.} Si no sabe escribirlo, aún
          no sabe qué tiene que automatizar.
    \item \textbf{Ponga el criterio en el programa.} Toda aritmética, toda
          comparación, toda conversión: en código, no en el prompt.
    \item \textbf{Valide la forma, no confíe en ella.} Un JSON que se parsea
          no es un JSON que cumple su esquema.
    \item \textbf{Defina la escala en el instrumento.} Y si no la hay, no
          publique el número.
    \item \textbf{Verifique por registro, no por bandera verde.} Un programa
          que compila o termina con 0 puede estar mintiendo: lea su salida.
    \item \textbf{Deje la frontera escrita.} Los ficheros con LLM se nombran;
          todo lo demás es código.
  \end{enumerate}
\end{frame}

\begin{frame}[plain]
  \vfill
  \begin{center}
    {\Large\bfseries\color{fondoOscuro} IA donde aporta.}\\[4pt]
    {\Large\bfseries\color{azulNoche} Código donde se exige.}\\[4pt]
    {\Large\bfseries\color{verdeTurquesa} Verificación en ambos.}\\[16pt]
    {\normalsize\color{grisTexto}
      La arquitectura no es una jaula para el modelo:\\
      es lo que le permite ser útil sin dejar de ser verificable.}
  \end{center}
  \vfill
\end{frame}

\end{document}
